\documentclass[11pt,a4paper]{article}
\usepackage[utf8]{inputenc}
\usepackage[T1]{fontenc}
\usepackage{lmodern}
\usepackage[margin=2cm,top=2.2cm,bottom=2.2cm]{geometry}
\usepackage{booktabs}
\usepackage{enumitem}
\usepackage{fancyhdr}
\usepackage{titlesec}
\usepackage{xcolor}
\usepackage{listings}
\usepackage{microtype}
\usepackage[colorlinks=true,linkcolor=blue!70!black,citecolor=blue!70!black,urlcolor=blue!70!black]{hyperref}

% Header and Footer styling
\pagestyle{fancy}
\fancyhf{}
\fancyhead[L]{\small\textsf{\textbf{CI \& Pre-commit Guide} \textbar\ Guardian Layer}}
\fancyhead[R]{\small\textsf{Page \thepage}}
\fancyfoot[L]{\small\textsf{DevOps \& Quality Assurance Documentation}}
\fancyfoot[R]{\small\textsf{Version 1.0}}
\renewcommand{\headrulewidth}{0.5pt}
\renewcommand{\footrulewidth}{0.5pt}

% Section formatting
\titleformat{\section}
  {\normalfont\sffamily\Large\bfseries\color{blue!80!black}}
  {\thesection}{0.8em}{}
\titleformat{\subsection}
  {\normalfont\sffamily\large\bfseries\color{blue!70!black}}
  {\thesubsection}{0.6em}{}

\titlespacing*{\section}{0pt}{10pt}{5pt}
\titlespacing*{\subsection}{0pt}{8pt}{4pt}
\setlength{\parskip}{3pt}
\setlength{\parindent}{0pt}

% Code listing styling
\definecolor{codegray}{rgb}{0.96,0.96,0.96}
\definecolor{codeframe}{rgb}{0.85,0.85,0.85}
\lstset{
  backgroundcolor=\color{codegray},
  basicstyle=\ttfamily\footnotesize,
  breaklines=true,
  frame=single,
  rulecolor=\color{codeframe},
  keepspaces=true
}

\title{\sffamily\bfseries\huge CI and Pre-Commit Guide\\[0.3em]
\Large Quality Gates, Code Standards, and Automated CI Pipeline}
\author{\textsf{The Guardian Layer Engineering Team}}
\date{\textsf{Version 1.0 \textbar\ November 2026}}

\begin{document}

\maketitle
\thispagestyle{fancy}

\section{Purpose}
The Guardian Layer codebase enforces automated quality gates at both local developer environments (via git pre-commit hooks) and pull-request integration perimeters (via GitHub Actions workflows). These verification layers ensure zero regressions in formatting, static analysis, type checking, unit/integration testing, and commit conventions before any code merges into the mainline branch.

\section{Pre-Commit Hooks}
Local commit validation is orchestrated via \texttt{.pre-commit-config.yaml}. The hook suite executes automatically during \texttt{git commit}:

\begin{itemize}[leftmargin=*,itemsep=2pt]
  \item \textbf{Ruff format:} Enforces PEP 8 code formatting standards across all Python sources.
  \item \textbf{Ruff lint:} Evaluates Python code quality rules, catching dead code, unused imports, and anti-patterns.
  \item \textbf{Tests:} Executes the fast pytest test suite, preventing regression commits.
  \item \textbf{Commitizen:} Validates commit message adherence to Conventional Commits conventions (e.g., \texttt{feat:}, \texttt{fix:}, \texttt{chore:}).
  \item \textbf{Branch check:} Verifies feature and bugfix branch naming standards.
\end{itemize}

\subsection{Installation and Execution}
To install hooks locally into the git repository:
\begin{lstlisting}
uv run pre-commit install
uv run pre-commit install --hook-type commit-msg
\end{lstlisting}

To manually trigger the full pre-commit evaluation against all files:
\begin{lstlisting}
uv run pre-commit run --all-files
\end{lstlisting}

\section{CI Pipeline Architecture}
The central continuous integration pipeline is defined in \texttt{.github/workflows/ci.yml}. Triggered on every Pull Request targeting the primary branch, the workflow executes the following sequential stages:

\begin{enumerate}[leftmargin=*,itemsep=2pt]
  \item \textbf{Repository Checkout:} Clones repository source with full commit history.
  \item \textbf{Runtime Provisioning:} Sets up Python 3.12 and resolves frozen lockfile dependencies via \texttt{uv sync}.
  \item \textbf{Linter Validation:} Runs \texttt{ruff check} against all codebase modules.
  \item \textbf{Type Safety Evaluation:} Runs \texttt{basedpyright} type analysis, enforcing strict typing contracts.
  \item \textbf{Test Suite Execution:} Runs all unit and integration tests via \texttt{pytest}.
  \item \textbf{Container Build Check:} Discovers Dockerfiles within repository roots and verifies successful image builds.
\end{enumerate}

\noindent\textit{Failure policy:} If any linter rule, type check, pytest assertion, or container build fails, the CI job fails immediately and blocks branch merging.

\section{Docker Compose Environment}
The \texttt{docker-compose.yml} configuration provides multi-service local infrastructure:
\begin{itemize}[leftmargin=*,itemsep=2pt]
  \item \textbf{Kafka Brokers:} 3-node Apache Kafka cluster using image \texttt{apache/kafka:3.8.0} with KRaft consensus.
  \item \textbf{Relational Database:} PostgreSQL 16 container for persistent telemetry state and anomaly verdict tables.
  \item \textbf{Audit Logging:} Elasticsearch 8.15 container storing operational telemetry and forensic logs.
\end{itemize}

\section{Pre-PR Submission Checklist}
Before opening a Pull Request, engineers must execute the following local verification sequence:
\begin{lstlisting}
uv run pre-commit run --all-files
git diff --check
git status
\end{lstlisting}
Ensure all local hooks pass and all changes are cleanly staged before requesting team review.

\end{document}

